\input{header.tex}
\input{newcommands.tex}
\singlespacing

%Skizze zu Aufgabe (1): Kugelstoßen, 1 m entspricht 0,5 cm.
%Aufgabenbild: \skizzeKugelstoss     Lösungsbild mit Koordinatensystem: {\mitKStrue \skizzeKugelstoss}
\newcommand{\skizzeKugelstoss}{%
\begin{tikzpicture}[x=0.5cm, y=0.5cm]
	\ifmitKS \skizzenraster{-1}{11}{0}{4}{1}{2,4,6,8,10}{1,2,3,4}
	\else \fill[orange!25] (-2,-0.45) rectangle (11.2,0); \fi
	\draw[thick, orange!50!black] (-2,0) -- (11.2,0);
	\emojifigur[xscale=-1]{(-0.85,0)}{person playing handball}{1.25cm}
	\ifmitKS
		\draw[akzentblau, thick, domain=0:10, samples=100] plot (\x,{-0.1*(\x-4)^2+3.6});
		\skizzenpunkt{(0,2)}{above left}{$(0\,|\,2)$}
		\skizzenpunkt{(4,3.6)}{above}{$S(4\,|\,3{,}6)$}
		\skizzenpunkt{(10,0)}{above right}{$N(10\,|\,0)$}
	\else
		\draw[dashed, thick, gray!60!black, domain=0:10, samples=100] plot (\x,{-0.1*(\x-4)^2+3.6});
		\shade[ball color=gray] (6.5,2.975) circle (0.24);
		\draw[densely dotted, thick] (6.5,0) -- (6.5,2.72) node[midway, right, font=\scriptsize, inner sep=2pt] {$h(x)$};
		\masslinie{(0,-1.1)}{(6.5,-1.1)}{\scriptsize $x$}
	\fi
\end{tikzpicture}}

\begin{document}

	\begin{center}\textbf{\Large{Sachaufgabe mit Skizze}}\end{center}
	\bigskip

	\begin{enumerate}[(1)]
		\item
			\begin{minipage}[t]{0.56\linewidth}
				Die Flugbahn einer Kugel beim Kugelstoßen lässt sich durch die Funktion $h$ mit $h(x)=-0{,}1(x-4)^2+3{,}6$ beschreiben. Dabei ist $x$ die waagerechte Entfernung vom Abstoßpunkt in Metern und $h(x)$ die Höhe der Kugel in Metern.
				\begin{enumerate}[(a)]
					\item Berechne die Höhe, in der die Kugel die Hand verlässt.
					\item Gib die maximale Höhe der Kugel und die Entfernung an, in der sie erreicht wird.
					\item Berechne die Stoßweite.
				\end{enumerate}
			\end{minipage}\hfill
			\begin{minipage}[t]{0.42\linewidth}
				\vspace{0pt}\centering
				\skizzeKugelstoss
			\end{minipage}
	\end{enumerate}

	\newpage

	\begin{center}\textbf{\Large{Lösungen}}\end{center}
	\medskip
	\begin{scriptsize}\begin{enumerate}[(1)]
		\item
			\begin{minipage}[t]{0.60\linewidth}
				\vspace{0pt}\raggedright
				\begin{enumerate}[(a)]
					\item $h(0)=-0{,}1\cdot(-4)^2+3{,}6=2$:\\ $\boxed{\text{Die Kugel verlässt die Hand in } \unit[2]{m} \text{ Höhe.}}$
					\item Scheitelpunkt $S(4\,|\,3{,}6)$:\\ $\boxed{\text{maximale Höhe } \unit[3{,}6]{m} \text{ in } \unit[4]{m} \text{ Entfernung}}$
					\item $-0{,}1(x-4)^2+3{,}6=0 \;\Rightarrow\; (x-4)^2=36 \;\Rightarrow\; x_1=-2,\; x_2=10$; \;$x_1$ entfällt.\\
						$\boxed{\text{Die Stoßweite beträgt } \unit[10]{m}.}$
				\end{enumerate}
			\end{minipage}\hfill
			\begin{minipage}[t]{0.38\linewidth}
				\vspace{0pt}\centering
				\resizebox{\linewidth}{!}{{\mitKStrue \skizzeKugelstoss}}
			\end{minipage}
	\end{enumerate}\end{scriptsize}

\end{document}
